\section{Conclusion}
\label{sec:conclude}

Link-variable smoothing is a crucial ingredient in constructing
gluonic operators which have dramatically reduced mixings with
the high frequency modes of the theory.  Link smearing is also
playing an increasingly important role in the construction of
improved lattice actions.  The lack of differentiability with respect
to the underlying link variables of standard smearing schemes
prevents the use of efficient Monte Carlo updating methods based
on molecular dynamics evolution.  

A link smearing method which circumvents these problems was proposed and
tested in this paper. The link smearing method is analytic everywhere in
the finite complex plane and utilizes the exponential function in such a
manner to remain within $SU(3)$, eliminating the need for any projection
back into the group.  Because of this construction, the algorithm is
also useful for any Lie group. An efficient implementation of this
smearing scheme, as well as the recursive computation of the force term
describing the change of the action in response to a variation of the link
variables, was described.  Although the force field is related to the
underlying link variables in a very complicated manner, each step
in the recursive computation of the force field is actually straightforward,
allowing a natural and efficient design in software.  The smeared mean
plaquette and the effective energy associated with the static 
quark-antiquark potential were used to show that no degradation in
effectiveness is observed as compared to link smearing
methods currently in use, although an increased sensitivity to the
smearing parameter was found.

